% answers_short.tex —— 2014 年·天津卷·文科 · 简短答案速查
% 由 tools/make_short_answers.py 自 answers.tex 的 \gksolhead/\gkanswer 提取，
% 并与 2014/2014_数学真题_天津文_答案要点.md 逐题比对后生成，勿手改（改源 MD 后重跑）。

\begin{gknotice}
\textbf{使用说明}：本节只给最终答案便于快速核对，完整推导见解析版。答案为编写者独立推导，\textbf{非官方标准答案}。
\end{gknotice}

\gkpart{选择题}
\begin{center}\begin{tabular}{ccccccccc}
\toprule
题号 & 1 & 2 & 3 & 4 & 5 & 6 & 7 & 8 \\
\midrule
答案 & A & B & B & C & D & A & D & C \\
\bottomrule
\end{tabular}\end{center}

\gkpart{填空题}
\begin{itemize}[leftmargin=2.6em]
  \item[\textbf{9.}] $60$。总份数 $20$，每份 $15$ 人，一年级 $4$ 份即 $60$ 人。
  \item[\textbf{10.}] $\dfrac{20\pi}{3}$。几何体为半径 $1$、高 $4$ 的圆柱与半径 $2$、高 $2$ 的圆锥同轴叠放，$4\pi+\dfrac{8\pi}{3}=\dfrac{20\pi}{3}$。
  \item[\textbf{11.}] $-4$。依次得 $S=-8$（$n=3$）、$S=-4$（$n=2$），此时 $n=1\leqslant 1$，输出 $-4$。
  \item[\textbf{12.}] $(-\infty,0)$。$f(x)=2\lg|x|$ 在 $x<0$ 时随 $x$ 增大而减小，在 $x>0$ 时递增。
  \item[\textbf{13.}] $2$。由 $\overrightarrow{AE}=\boldsymbol u+\dfrac13\boldsymbol v$、$\overrightarrow{AF}=\boldsymbol v+\dfrac1\lambda\boldsymbol u$ 及 $\boldsymbol u\cdot\boldsymbol v=-2$ 得 $-\dfrac23+\dfrac{10}{3\lambda}=1$。
  \item[\textbf{14.}] $(1,2)$。$x>0$ 的零点数在 $0<a<2$ 时为 $2$、$a\geqslant 2$ 时为 $1$；$x<0$ 的零点数在 $0<a<1$ 时为 $4$、$a=1$ 时为 $3$、$a>1$ 时为 $2$。合计恰为 $4$ 当且仅当 $1<a<2$。
\end{itemize}

\gkpart{解答题}
\begin{itemize}[leftmargin=2.6em]
  \item[\textbf{15.}] (1) $AB$，$AC$，$AX$，$AY$，$AZ$，$BC$，$BX$，$BY$，$BZ$，$CX$，$CY$，$CZ$，$XY$，$XZ$，$YZ$，共 $\mathrm C_6^2=15$ 种；(2) $\dfrac25$。符合条件的结果为 $AY$，$AZ$，$BX$，$BZ$，$CX$，$CY$ 共 $6$ 种，$P(M)=\dfrac{6}{15}=\dfrac25$
  \item[\textbf{16.}] (1) $\dfrac{\sqrt6}{4}$。由 $b=\sqrt6c$ 得 $a=2c$，再用余弦定理；(2) $\dfrac{\sqrt{15}-\sqrt3}{8}$。由 $\cos2A=-\dfrac14$、$\sin2A=\dfrac{\sqrt{15}}{4}$ 及和角公式
  \item[\textbf{17.}] (1) 见解析版；(2) ① 由 $AB\perp BD$、$PO\perp AD$、$BO\perp AD$ 得 $AD\perp$ 平面 $POB$，$\angle POB=60^\circ$，算得 $PB\perp BO$，从而 $PB\perp$ 平面 $ABCD$，故平面 $PBC\perp$ 平面 $ABCD$；② $\dfrac{2\sqrt{11}}{11}$。以 $B$ 为原点建系，$\overrightarrow{EF}=\left(0,-\sqrt2,\dfrac{\sqrt3}{2}\right)$，平面 $PBC$ 的法向量 $(1,1,0)$
  \item[\textbf{18.}] (1) $\dfrac{\sqrt2}{2}$。由 $a^2+b^2=3c^2$ 与 $b^2=a^2-c^2$ 得 $a^2=2c^2$；(2) $\dfrac{x^2}{6}+\dfrac{y^2}{3}=1$。由 $\overrightarrow{F_1P}\cdot\overrightarrow{F_1B}=0$ 得 $P\left(-\dfrac{4c}{3},\dfrac c3\right)$，再由 $|C'F_2|^2=r^2+|MF_2|^2$ 得 $c^2=3$
  \item[\textbf{19.}] (1) 递减区间 $(-\infty,0)$ 与 $\left(\dfrac1a,+\infty\right)$，递增区间 $\left(0,\dfrac1a\right)$；极小值 $f(0)=0$，极大值 $f\left(\dfrac1a\right)=\dfrac{1}{3a^2}$；(2) $\left[\dfrac34,\dfrac32\right]$。须 $f$ 在 $(2,+\infty)$ 上恒负，即 $a\geqslant\dfrac34$；又 $f$ 在 $(1,+\infty)$ 上值域的上端须不小于 $0$，即 $a\leqslant\dfrac32$
  \item[\textbf{20.}] (1) $A=\{0,1,2,3,4,5,6,7\}$。$q=2$、$n=3$ 时 $x=x_1+2x_2+4x_3$，三系数各取 $0$ 或 $1$，共 $2^3=8$ 个互异值；(2) 见解析版
\end{itemize}
